% =====================================================================
\chapter{Derivadas matriciales: capa lineal y softmax}
\label{tema:3}
% =====================================================================

Casi todo el cómputo del entrenamiento de un modelo de lenguaje consiste en derivar dos tipos de operación: capas lineales y la combinación de softmax con entropía cruzada en la salida. Este tema define el gradiente de una función escalar respecto de una matriz, obtiene los gradientes de la capa lineal, individualmente y por lotes, y deduce el resultado más usado del curso: el gradiente de la entropía cruzada respecto de los logits es $p-y$. Se sigue~\cite[§5.3--§5.5]{mml}, la tarea~2 de~\cite{cs224n} y, como tabla de consulta, \cite{cookbook}.

\section{Gradiente respecto de una matriz}

\begin{definicion}[Gradiente respecto de una matriz]\label{def:grad-matriz}
Sea $L\colon\R^{m\times n}\to\R$. Su \emph{gradiente} en $W$ es la matriz $\nabla_W L=\partial L/\partial W\in\R^{m\times n}$ de entradas
\[
\Bigl(\frac{\partial L}{\partial W}\Bigr)_{ij}=\frac{\partial L}{\partial W_{ij}}.
\]
Equivale a identificar $\R^{m\times n}$ con $\R^{mn}$ y reordenar el gradiente con la forma de $W$.
\end{definicion}

\begin{proposicion}[Aproximación lineal]\label{prop:frobenius}
Si $L$ es diferenciable en $W$, para toda perturbación $E\in\R^{m\times n}$,
\[
L(W+E)=L(W)+\Bigl\langle\frac{\partial L}{\partial W},E\Bigr\rangle+o(\norm{E}),\qquad
\langle A,B\rangle=\sum_{i,j}A_{ij}B_{ij}=\tr(A\T B).
\]
\end{proposicion}

La regla práctica que se deriva es que \emph{el gradiente de un escalar respecto de un parámetro tiene la misma forma que el parámetro}. Comprobar las formas es la primera verificación de cualquier fórmula de este tema.

\section{La capa lineal}

\begin{proposicion}[Gradientes de una capa lineal]\label{prop:capa}
Sean $W\in\R^{m\times n}$, $x\in\R^{n}$, $b\in\R^{m}$, $y=Wx+b$ y $L=\ell(y)$ con $\ell\colon\R^{m}\to\R$ diferenciable. Sea $\delta=\nabla_y\ell\in\R^{m}$ la \emph{señal de error} que llega a la salida. Entonces
\[
\frac{\partial L}{\partial W}=\delta\,x\T,\qquad
\frac{\partial L}{\partial b}=\delta,\qquad
\frac{\partial L}{\partial x}=W\T\delta.
\]
\end{proposicion}

Las formas encajan: $\delta x\T$ es $(m\times 1)(1\times n)=m\times n$, como $W$, y $W\T\delta$ es $(n\times m)(m\times 1)=n\times 1$, como $x$. El vector $W\T\delta$ es la señal de error que la capa transmite a la capa anterior.

\addtocounter{definicion}{1}%

\begin{corolario}[Capa lineal por lotes]\label{cor:lotes}
Sean $X\in\R^{B\times n}$ una matriz cuyas filas son $B$ ejemplos, $W\in\R^{n\times m}$, $Z=XW\in\R^{B\times m}$ y $L=\ell(Z)$ con $\Delta=\partial L/\partial Z\in\R^{B\times m}$. Entonces
\[
\frac{\partial L}{\partial W}=X\T\Delta,\qquad\frac{\partial L}{\partial X}=\Delta\,W\T.
\]
\end{corolario}

El \cref{cor:lotes} es la \cref{prop:capa} escrita con filas en lugar de columnas y sumada sobre los ejemplos del lote: $X\T\Delta=\sum_b x_b\,\delta_b\T$. Es la convención de las bibliotecas de aprendizaje profundo.

\addtocounter{definicion}{1}%

\section{El jacobiano de la softmax}

Se recuerda del módulo~1 que, para $z\in\R^{K}$, $p=\softmax(z)$ tiene componentes $p_i=e^{z_i}/S$ con $S=\sum_k e^{z_k}$.

\begin{proposicion}[Jacobiano de la softmax]\label{prop:jac-softmax}
\[
\frac{\partial p_i}{\partial z_j}=p_i\bigl(\mathbb{I}[i=j]-p_j\bigr),\qquad\text{es decir,}\qquad J_{\softmax}(z)=\diag(p)-pp\T.
\]
Esta matriz es simétrica y cumple $J_{\softmax}(z)\,\uno=0$.
\end{proposicion}

La igualdad $J_{\softmax}\uno=0$ es la versión infinitesimal de la invariancia por traslación de la softmax (módulo~1, tema~4): sumar la misma constante a todos los logits no cambia las probabilidades.

\section{Entropía cruzada y el gradiente \texorpdfstring{$p-y$}{p − y}}

Sea $c$ el índice del token correcto e $y=e_c$ su vector \emph{one-hot}. La pérdida de entropía cruzada es $L(z)=-\ln p_c$.

\begin{teorema}[Gradiente de la entropía cruzada]\label{teo:p-y}
\[
\nabla_z L=p-y.
\]
\end{teorema}

\addtocounter{definicion}{1}%

\begin{observacion}[Interpretación]
Un paso de descenso $z\leftarrow z-\eta(p-y)$ aumenta el logit correcto en $\eta(1-p_c)$ y disminuye cada logit incorrecto en $\eta p_i$: se penaliza más a los competidores más probables. El gradiente es pequeño cuando el modelo acierta con confianza ($p_c\approx 1$) y grande cuando se equivoca. Como la expresión $p-y$ no contiene cocientes por $p_c$, las bibliotecas implementan la softmax y la entropía cruzada como una sola operación, más estable que componer las dos (módulo~11).
\end{observacion}

\begin{corolario}[Capa de salida de un LLM]\label{cor:salida}
Sean $h\in\R^{d}$ el estado oculto de la última posición, $W\in\R^{K\times d}$ la matriz de salida, $z=Wh$ y $L=-\ln\softmax(z)_c$. Entonces
\[
\frac{\partial L}{\partial W}=(p-y)\,h\T,\qquad\frac{\partial L}{\partial h}=W\T(p-y).
\]
\end{corolario}

Con un vocabulario de $K=10^{5}$ tokens, $W$ tiene $10^{5}$ filas; la fila $i$ recibe la corrección $(p_i-y_i)\,h$, de modo que todas las filas se actualizan en cada paso, no solo la del token correcto.

